\chapter{Fundamentação Teórica}
\label{chap:fundamentacao}

A Fundamentação Teórica estabelece a base conceitual e matemática necessária para o pleno entendimento da proposta metodológica e dos resultados do trabalho. Nesta seção, o autor apresenta os princípios, algoritmos, modelos formais e tecnologias sobre os quais a pesquisa é construída.

Diferente dos Trabalhos Relacionados, que analisam soluções concorrentes ou correlatas, a Fundamentação Teórica foca nos \textbf{conceitos consolidados} da literatura científica, citando autores de referência no domínio \cite{Ding2020AdversarialAttacks}.

Para manter a legibilidade, o capítulo deve ser dividido em seções temáticas coerentes. Quando o trabalho envolver conceitos matemáticos ou formulações analíticas, deve-se empregar o ambiente formal de equações:

\begin{equation}
    x = \frac{-b \pm \sqrt{b^2 - 4ac}}{2a}
    \label{eq:quadratic}
\end{equation}

Toda equação deve ser numerada e explicitamente discutida no texto (ex.: \enquote{conforme a Equação~\ref{eq:quadratic}}). Quando pertinente, a inclusão de trechos de código-fonte ilustrativos pode ser feita utilizando o pacote \texttt{listings}, como demonstrado no Código~\ref{lst:exemplocodigo}.

\lstinputlisting[language=Python,
caption=Exemplo de listagem de código-fonte estruturado,
label=lst:exemplocodigo]{content/\contentlangdir/chapters/fundamentacao/example.py}
\hfill
\begin{minipage}[t]{.65\textwidth}
\fonteautor
\end{minipage}
